%
% 03_einfache_Herleitung.tex -- Das dritte Gesetz im Spezialfall der
% Kreisbahn, einmal elementar und einmal als Spezialfall des Integrals
%
% !TEX root = ../../buch.tex
% !TEX encoding = UTF-8
%
\section{Der Spezialfall der Kreisbahn
\label{kepler:section:kreisbahn}}

Wir wagen uns noch nicht direkt an die allgemeine Herleitung für Ellipsen.
Zuerst betrachten wir den Spezialfall einer perfekten Kreisbahn.
Dieser Fall lässt sich mit elementarer Mechanik behandeln.
Er dient uns später als Kontrolle für das allgemeine Resultat.
Auf einer Kreisbahn ist der Abstand zur Sonne konstant.
Die grosse Halbachse $a$ stimmt dann mit dem Radius $r$ überein.

\subsection{Herleitung aus dem Kräftegleichgewicht
\label{kepler:subsection:kraeftegleichgewicht}}
\kopfrechts{Der Spezialfall der Kreisbahn}%

Ein Planet der Masse $m$ soll auf einer Kreisbahn mit Radius $r$ bleiben.
Dazu muss die Gravitationskraft der Sonne genau die nötige Zentripetalkraft aufbringen.
\index{Zentripetalkraft}%
Mit der Winkelgeschwindigkeit $\omega=2\pi/T$ lautet diese Bedingung
\index{Winkelgeschwindigkeit}%
\begin{equation}
G\,\frac{mM}{r^2}
=
m\,\omega^2 r
=
m\,\frac{4\pi^2}{T^2}\,r.
\label{kepler:eqn:gleichgewicht}
\end{equation}
Die Planetenmasse $m$ kürzt sich auf beiden Seiten heraus.
Dann bringen wir die Bahndaten $T$ und $r$ auf eine Seite.
Das ergibt
\begin{equation}
\frac{T^2}{r^3}
=
\frac{4\pi^2}{GM}
=
\text{konstant}.
\label{kepler:eqn:kreisbahn}
\end{equation}
Das ist das dritte keplersche Gesetz für Kreisbahnen.
Der Quotient $T^2/r^3$ hängt nur von der Masse $M$ des Zentralgestirns ab.
\index{Zentralgestirn}%
Vom umlaufenden Körper hängt er nicht ab.

Diese Rechnung ist einfach.
Für reale Planeten beweist sie aber wenig, denn deren Bahnen sind keine Kreise.
Kepler fand seine Gesetze an den Bahndaten des Mars.
Schon dessen Bahn hat die Exzentrizität $\varepsilon\approx 0.09$.
Bei Kometen kann die Exzentrizität sogar nahe bei $1$ liegen.
Ein allgemeiner Beweis braucht deshalb die Umlaufzeit auf einer Ellipse.
Wir müssen also das Integral in \eqref{kepler:eqn:tintegral} auswerten.

\subsection{Die Kreisbahn als Spezialfall des Integrals
\label{kepler:subsection:integralspezialfall}}

Als Vorbereitung lösen wir das Integral in \eqref{kepler:eqn:tintegral} zuerst für die Kreisbahn.
Das ist gleichzeitig ein Konsistenztest.
Eine Kreisbahn hat die Exzentrizität $\varepsilon=0$.
Der Integrand wird dadurch konstant und das Integral trivial:
\[
\int_0^{2\pi}
\frac{1}{(1+0\cdot\cos(\vartheta))^2}
\,d\vartheta
=
\int_0^{2\pi} 1\,d\vartheta
=
2\pi.
\]
Aus \eqref{kepler:eqn:tintegral} wird damit
\[
T
=
\frac{m\,p^2}{L}\cdot 2\pi.
\]
Nun ersetzen wir den Drehimpuls mit der Beziehung \eqref{kepler:eqn:drehimpulsgeometrie} durch $L=m\sqrt{GMp\mathstrut}$.
Die Planetenmasse kürzt sich wieder heraus:
\[
T
=
\frac{m\,p^2}{m\sqrt{GMp\mathstrut}}\cdot 2\pi
=
\frac{2\pi\,p^{\frac{3}{2}}}{\!\sqrt{GM\mathstrut\rlap{$\phantom{p}$}}}.
\]
Auf einer Kreisbahn sind beide Halbachsen und der Halbparameter gleich dem Radius.
Es gilt also $p=r$.
Quadrieren liefert
\[
T^2
=
\frac{4\pi^2}{GM}\,r^3,
\]
in Übereinstimmung mit \eqref{kepler:eqn:kreisbahn}.
Der Integralansatz führt für $\varepsilon=0$ also auf dieselbe Konstante wie das Kräftegleichgewicht.
Offen bleibt der Fall $0<\varepsilon<1$.
Dort müssen wir zeigen: Mit der grossen Halbachse $a$ anstelle von $r$ gilt das Gesetz weiterhin.
